% Generated draft: requires mathematical and source review.
\documentclass[11pt]{article}
\usepackage[margin=25mm]{geometry}
\usepackage{fontspec}
\setmainfont{Latin Modern Roman}
\usepackage{amsmath,amssymb,mathtools,unicode-math}
\setmathfont{Latin Modern Math}
\usepackage{enumitem,needspace,xurl,hyperref}
\setlength{\parindent}{0pt}
\setlength{\parskip}{5pt}
\setlength{\emergencystretch}{3em}
\newcommand{\N}{\mathbb{N}}
\newcommand{\Z}{\mathbb{Z}}
\newcommand{\Q}{\mathbb{Q}}
\newcommand{\R}{\mathbb{R}}
\newcommand{\C}{\mathbb{C}}
\newcommand{\F}{\mathbb{F}}
\newcommand{\A}{\mathbb{A}}
\newcommand{\PP}{\mathbb{P}}
\newcommand{\E}{\mathbb{E}}
\newcommand{\eps}{\varepsilon}
\newcommand{\dd}{\,\mathrm{d}}
\begin{document}
\section*{UnsolvedMath ID: 10400036}
\textbf{Source problem code:} catalogue entry
\subsection{Definitions and mathematical statement}
An ordered oriented $n$-component string link is a proper embedding of $n$ intervals in $D^2\times[0,1]$, with fixed corresponding endpoints on the bottom and top disks. With the standard meridian, longitude and basing conventions, its Milnor numbers are defined by Magnus coefficients of longitudes in the appropriate nilpotent quotients: a meridian $m_j$ is expanded as $1+X_j$, and $\mu_{i_1\cdots i_r,j}$ is the coefficient of $X_{i_1}\cdots X_{i_r}$ in the $j$th longitude expansion.

The problem attributed to M. Polyak asks for a topological presentation of these invariants without assuming lower Milnor invariants vanish. Specifically, it requests a meaning for the derived object denoted $L_{12\cdots n-1}$ such that
\[
\operatorname{lk}(L_{12\cdots n-1},L_n)
\]
is well defined in an appropriate sense, and then an identification
\[
\operatorname{lk}(L_{12\cdots n-1},L_n)
=\mu_{12\cdots n-1,n}(L).
\]
Both well-definedness and equality are part of the task. The source does not reproduce the preceding construction of $L_{12\cdots n-1}$, its choices, closure/chains for linking, or what ``appropriate sense'' means. One cannot replace it by an arbitrary iterated intersection or impose lower-invariant vanishing to avoid the issue.

In the usual based string-link convention, integer Milnor invariants can be defined without the closed-link indeterminacies. The supplied summary's statement that indeterminacy is inherent must therefore be checked against the chosen object and basing: string-link invariants, closed-link $\overline\mu$ invariants, and ambiguity of this geometric linking construction are different issues. No modulus or quotient is guessed.

The extract then enters a separate section 2.7 on virtual knots, not an added hypothesis of Problem 2.14. It describes virtual-crossing diagrams modulo Reidemeister moves, arrow diagrams as oriented chords, weight systems on the arrow-diagram module modulo 6T and oriented FI relations, and an isomorphism with acyclic oriented Jacobi diagrams modulo oriented AS, IHX, STU and pictured extra relations. The skeleton $X$, Figure 9 and the extra pictures are not reproduced, so no full algebraic presentation or additional virtual-knot conjecture is reconstructed from that trailing context.
\subsection{Short English statement}
Give a geometric linking-number description of string-link Milnor invariants with no lower-invariant vanishing assumption. Define the source's derived linking object, prove its appropriate well-definedness and identify it with the specified Milnor number.
\subsection{Sources}
\begin{itemize}
\item Local UnsolvedMath registry, permanent ID 10400036, exact source code catalogue entry, Problem 2.14; source attribution M. Polyak, no proposer/year in metadata. Registry status partially solved, classification PARTIAL-PROGRESS. The registry title is truncated.
\item Ohtsuki, \emph{Problems on invariants of knots and 3-manifolds} (2002), Problem 2.14, PDF page 37, \url{https://msp.org/gtm/2002/04/gtm-2002-04-024s.pdf}. Supplied pdf-font-structure extraction accessed July 29, 2026; provenance distinguishes 231 bold numbered items. No PDF was accessed here.
\item The source cites [329] for string-link Milnor invariants. Its trailing virtual-knot section cites [203], [154] and [330], with Figure 9 for relations. The bibliography entries and figure are absent, so no author/title is invented for these labels.
\item The supplied summary mentions topological interpretation and indeterminacy but gives no precise later result; its closed-link/string-link scope requires checking. Provenance flags current primary-literature status and redistribution terms NEEDS\_REVIEW.
\end{itemize}
\end{document}
